\documentclass[zihao=5, a4paper]{ctexart}

% -----------------------------------------------------------
% 宏包引入
% -----------------------------------------------------------
\usepackage{amsmath, amssymb, amsthm} % 数学符号
\usepackage{geometry}                 % 页面设置
\usepackage{enumitem}                 % 列表优化
\usepackage{xcolor}                   % 颜色（用于提示填空）


\geometry{left=2.5cm, right=2.5cm, top=2.5cm, bottom=2.5cm}
\linespread{1.25} % 五号字体的行距微调

% -----------------------------------------------------------
% 自定义命令（分值）
% -----------------------------------------------------------
\newcommand{\fillhere}[1]{\textcolor{red}{\textbf{[#1]}}} 
\newcommand{\scoreB}[1]{\par\dotfill \textbf{(#1分)}\\}

\begin{document}

% ===========================================================
% 第一部分：原创声明
% ===========================================================
\section*{一、原创声明}

\begin{center}
\framebox[0.9\textwidth][c]{
\begin{minipage}{0.85\textwidth}
\vspace{1em}
\centerline{\large\textbf{声\quad 明}}
\vspace{1em}

本人郑重声明，所供题目为自己原创，本题目仅供给学而思培优使用，在被明确拒绝供题前或被接受后一年内，不会再将题目供至其他比赛中或在网上流传。本声明的法律结果由本人承担。

\vspace{2em}
\hfill 供题者：\underline{\makebox[3em][c]{戴天悦}} 

\hfill \underline{\makebox[3em][c]{2026}} 年 \underline{\makebox[2em][c]{2 }} 月 \underline{\makebox[2em][c]{ 23}} 日
\vspace{1em}
\end{minipage}
}
\end{center}

% ===========================================================
% 第二部分：个人信息
% ===========================================================
\section*{二、个人信息}

\begin{itemize}[label={}]
    \item \textbf{姓名}：戴天悦
    \item \textbf{身份证号}：440303200811075115
    \item \textbf{手机号}：13310830840
    \item \textbf{银行卡号}：6210817200020025035
    \item \textbf{开户行}：建设银行深圳鸿瑞支行
    \item \textbf{户主名}：戴波
\end{itemize}

\hrulefill

% ==========================================================================
% 第三部分：题目内容
% ==========================================================================
\section*{三、题目详情}

% 1. 建议位置-------------------------------------------------
\subsection*{1. 建议位置}
适合作为联赛的难度适中的第四题。
%-----------------------------------------------------------





% 2. 题干--------------------------------------------------
\subsection*{2. 题目}

学学和思思在玩一个游戏，首先给定一个正整数 $n$为四的倍数，学学有一个$n \times n$的方格表，并且这个方格表的上下两条边是粘起来的，左右两条边也是粘起来的（你可以把它想象成一个甜甜圈表面的网格）。游戏规则如下：
\begin{itemize}
    \item 思思需要在表格的每一个格子中填入 X、M、O 这三个字母中的任意一个。
    \item 填完后，学学会进行计分。学学会在网格中寻找"得分词"：\textbf{如果横向、纵向或 $45^\circ$ 对角线方向上，有连续的三个格子，它们连起来读恰好是拼写 X-M-O（或者是O-M-X）}，思思就能得 1 分。
    \item 由于表格是像甜甜圈一样连通的，跨越边界的相连格子也可以组成得分词。
\end{itemize}
请问，在这个游戏中，思思最多能获得多少分？

%-----------------------------------------------------------------


% 3. 解答---------------------------------------------------------------------
\subsection*{3. 解答}


所求得分的最大值是 $\frac{3}{2}n^2$。我们首先给出一个取等的构造：

将方格表从左到右按列依次填入字母 X、M、O、M，并以这 4 列为一个周期向右不断重复。在这种构造下，所有的 M 必定位于偶数列，而所有的 X 和 O 分别位于模 4 余 1 和余 3 的奇数列。经过验证，每一个 M 恰好能与相邻的 X 和 O 组成 3 个“得分词”（1 个横向，2 个对角线方向）。由于网格中共有 $\frac{n^2}{2}$ 个 M，总计恰好产生 $3 \times \frac{n^2}{2} = \frac{3}{2}n^2$ 个得分词。
\\

{\centering
\begin{tabular}{*{8}{c}}
X & M & O & M & X & M & O & M \\
X & M & O & M & X & M & O & M \\
X & M & O & M & X & M & O & M \\
X & M & O & M & X & M & O & M \\
X & M & O & M & X & M & O & M \\
X & M & O & M & X & M & O & M \\
X & M & O & M & X & M & O & M \\
X & M & O & M & X & M & O & M \\
\end{tabular}\par}
\scoreB{10}


我们将 $n \times n$ 的环面方格表等价视为 $\mathbb{Z}_n \times \mathbb{Z}_n$ 上的一个无向图 $G=(V, E)$。其中，顶点集 $V$ 代表方格表中的 $n^2$ 个格子。对于任意两个格子 $u, v \in V$，若它们在横向、纵向或 $45^\circ$ 对角线方向上相邻，则它们之间存在一条边。图中总边数为：
\[
|E| = \frac{8 \times n^2}{2} = 4n^2
\]

定义顶点赋值函数 $f: V \to \{1, -1\}$ 如下：
\begin{itemize}
    \item 若格子 $v$ 填入字母 M，则 $f(v) = 1$；
    \item 若格子 $v$ 填入字母 X 或 O，则 $f(v) = -1$。
\end{itemize}

\vspace{1em}

\large{\textbf{引理：}} \normalsize 设 $\#(XM)$ 为网格中连接 X 与 M 的边数，$\#(MO)$ 为连接 M 与 O 的边数。在任意一种填涂方案下，必然满足：
\[
\#(XM) + \#(MO) \le 3n^2
\]

\scoreB{20}

\textbf {\large{引理的证明：}} \normalsize

令 $\mathcal{X} = \#(XM) + \#(MO)$。根据赋值函数 $f(v)$，只有当一条边连接 M 与非 M（即 X 或 O）时，$(f(u) - f(v))^2 = (1 - (-1))^2 = 4$；否则 $(f(u) - f(v))^2 = 0$。因此：
\[
4\mathcal{X} = \sum_{(u,v) \in E} (f(u) - f(v))^2 = \sum_{(u,v) \in E} (2 - 2f(u)f(v))
\]
代入 $|E| = 4n^2$，化简得：
\[
4\mathcal{X} = 8n^2 - 2\sum_{(u,v) \in E} f(u)f(v)
\]

我们将网格划分为 $n^2$ 个 $2 \times 2$ 的小正方形(这里划分的意思是，环面中有$n^2$个 $2 \times 2$ 的小正方形)。对于任意一个正方形，设其四个顶点的赋值按位置分为 $a, b, c, d$。构造两个非负项：
\[
S_1 = (a + b + c + d)^2 \ge 0, \quad S_2 = (a - b - c + d)^2 \ge 0
\]
设 $P_H, P_V, P_D$ 分别为全局所有水平、垂直、对角线边的 $f(u)f(v)$ 乘积之和。对所有 $n^2$ 个正方形进行全局求和（注意到顶点被共享 4 次，水平/垂直边被共享 2 次，对角线边不共享）：
\[
\sum S_1 = 4n^2 + 4P_H + 4P_V + 2P_D \ge 0
\]
\[
\sum S_2 = 4n^2 - 4P_H - 4P_V + 2P_D \ge 0
\]
将两式按 $3 : 1$ 比例相加：
\[
\frac{3}{4} \sum S_1 + \frac{1}{4} \sum S_2 = 4n^2 + 2(P_H + P_V + P_D) \ge 0
\]
由于 $P_H + P_V + P_D = \sum_{(u,v) \in E} f(u)f(v)$，故得到：
\[
2n^2 + \sum_{(u,v) \in E} f(u)f(v) \ge 0 \implies \sum_{(u,v) \in E} f(u)f(v) \ge -2n^2
\]
代回 $\mathcal{X}$ 的等式中：
\[
4\mathcal{X} \le 8n^2 - 2(-2n^2) = 12n^2 \implies \mathcal{X} \le 3n^2
\]
引理得证。
\scoreB{40}

\vspace{1em}

根据题意，无论得分词是 X-M-O 还是 O-M-X，字母 M 必须严格位于中间。这意味着每一个“得分词”必然恰好对应两条共线的割边：一条 X-M 边和一条 M-O 边。

由于在同一条直线上，一对固定的相邻 X-M 边和 M-O 边只能组成一个确定的得分词，这意味着“得分词”的总数 $N$ 最多只能是 $\mathcal{X}$ 的一半：
\[
N \le \frac{1}{2} \mathcal{X}
\]
将引理 $\mathcal{X} \le 3n^2$ 代入，得到：
\[
N \le \frac{1}{2} (3n^2) = \frac{3}{2}n^2
\]
结合第一步的构造可知，该等号是可以取到的。故思思最多能获得 $\frac{3}{2}n^2$ 分。

\scoreB{50}
%-------------------------------------------------------------------------






% 4. 评分标准和附注-------------------------------------------------------
\subsection*{4. 评分标准}
\begin{itemize}
    \item 正确构造一个得分为 $\frac{3}{2}n^2$ 的方案：得 10 分。
    \item 完整声明出引理的内容，或者等价的重要中间结论：得 20 分。
    \item 完整给出引理的证明：得 40 分。引理证明的细节部分有小错误但不影响核心思路的，酌情扣 10 分。
    \item 正确利用引理证明最终结论：得 50 分。
    \item 任何其他可行的正确解法都应该获得满分。
\end{itemize}

\subsection*{5. 附注}
本做法可以推广到一般的没有重复字母的单词，比如ABCD，对于长度等于s的单词，我们可以证明最多只会出现$3/(s-1)n^2$个得分词，当n是$2s-2$的倍数时可以达到这个上界。
证明的办法是把字母分成在奇数位的字母和偶数位的字母，然后用相同的半不变量给出类似的引理，为了增加题目难度，可以把X-M-O换成别的单词。
但是作为联赛题，为避免细节上的繁琐，题干中采用了三个字母的版本。
\end{document}